% !TEX root = ../article.tex
% =============================================================================
% 3.3 Introducción
% =============================================================================
% Guía: contexto del problema y razones para abordarlo; usuarios o procesos
% que se beneficiarían; objetivo general, objetivos específicos y alcance;
% resultado esperado y forma de comprobar si se ha alcanzado.
%
% Primer avance (asesoría 1, 6-13 nov 2026): problema seleccionado, contexto
% y evidencia inicial, objetivos y resultados esperados.
% =============================================================================

\section{Introducción}\label{sec:introduction}

% Cifras y afirmaciones verificadas en los resúmenes de cada fuente (Crossref
% y PubMed) el 8 de noviembre de 2026.
\subsection{Contexto y motivación}\label{sec:context}
Las enfermedades pulmonares están entre las principales causas de muerte en el
mundo. En 2023, tres de las diez primeras afectaron al pulmón: la enfermedad
pulmonar obstructiva crónica (EPOC), tercera con unos 3,6 millones de muertes;
las infecciones de las vías respiratorias inferiores, cuarta con 2,7 millones;
y el cáncer de tráquea, bronquios y pulmón, séptima con 1,8 millones
\parencite{oms2026}. El cáncer de pulmón fue además, en 2024, el cáncer
diagnosticado con más frecuencia, con casi 2,6 millones de casos nuevos, y la
primera causa de muerte por cáncer, con unos 1,9 millones de fallecimientos
\parencite{sung2026}.

La tomografía computarizada (TC) permite examinar el pulmón en tres
dimensiones y medirlo. La TC cuantitativa traduce la imagen en medidas
numéricas automáticas, y una revisión sistemática de 185 estudios encontró que
estas medidas aportan información pronóstica en las enfermedades pulmonares
intersticiales \parencite{dixon2025}. Para obtenerlas, primero hay que separar
el pulmón del resto del tórax; esa delimitación, llamada segmentación, es un
paso crucial del análisis de imagen \parencite{hofmanninger2020}. En la
fibrosis pulmonar idiopática, por ejemplo, un modelo de Deep Learning segmentó
con éxito el 97,8~\% de las TC, de distintos fabricantes y con espesores de
corte de 0,5 a 5~mm. En una cohorte de 446 pacientes, el volumen pulmonar
obtenido se correlacionó con la capacidad vital forzada ($r = 0{,}82$), y un
volumen menor se asoció con una menor supervivencia libre de progresión a dos
años \parencite{thillai2024}.

Estas medidas, sin embargo, todavía no se usan de forma rutinaria en la
práctica clínica ni en los ensayos \parencite{thillai2024}, y no están
estandarizadas: la heterogeneidad de los métodos impidió comparar entre sí los
estudios de la revisión o combinarlos en un metaanálisis \parencite{dixon2025}.
Los métodos automáticos clásicos, basados en umbrales, regiones o formas, no
resuelven el problema, porque casi todos funcionan bien solo cuando el pulmón
presenta poca o ninguna patología; derrames pleurales, consolidaciones y masas
provocan segmentaciones inexactas \parencite{mansoor2015}. Delimitar el pulmón
a mano tampoco es viable en la rutina: en el conjunto de \textcite{ma2021},
delimitar el pulmón y la infección en una TC de 250 cortes tomó en promedio
unos 400 minutos. Además, reunir grandes conjuntos de imágenes anotadas en una
sola institución es poco práctico, sobre todo cuando los radiólogos están
dedicados a la atención \parencite{ma2021}.

El aprendizaje profundo ha cambiado este panorama. Con 70 casos de COVID-19
anotados, \textcite{ma2021} obtienen un coeficiente de Dice promedio de 97,3~\%
y 97,7~\% para el pulmón izquierdo y el derecho, aunque solo de 67,3~\% para la
infección. \textcite{hofmanninger2020} muestran que, en casos difíciles, la
precisión depende más de la diversidad de los datos de entrenamiento que de la
arquitectura: una U-Net entrenada con datos clínicos de rutina de varias
enfermedades alcanza un Dice de 0,98, frente a 0,94 de los métodos de
referencia. TotalSegmentator, un modelo nnU-Net entrenado con 1204 estudios de
TC de rutina, segmenta 104 estructuras anatómicas, con usos como la volumetría de órganos y la planificación de
radioterapia \parencite{wasserthal2023}. Persisten, no obstante, dos
dificultades: los métodos son difíciles de comparar porque se desarrollan con
datos, configuraciones y métricas distintos \parencite{ma2021}, y la
variabilidad propia del entrenamiento pone en duda la reproducibilidad de sus
resultados \parencite{renard2020}.

Thoraxis aborda la segmentación pulmonar en TC por cuatro razones:

\begin{enumerate}
    \item \textbf{Relevancia:} la segmentación del pulmón es el paso previo
          para cuantificar enfermedades con una carga sanitaria alta.
    \item \textbf{Viabilidad:} existen conjuntos de datos públicos con máscaras
          de referencia y modelos publicados que sirven como referencia, lo
          que permite evaluar la propuesta en condiciones comparables.
    \item \textbf{Brecha:} la robustez ante pulmones con patología y una
          evaluación reproducible y verificable siguen siendo problemas
          abiertos.
    \item \textbf{Proyección:} el pulmón es el primer caso; la misma
          infraestructura de entrenamiento, evaluación y despliegue puede
          extenderse a otras estructuras del tórax.
\end{enumerate}

Thoraxis es una herramienta de investigación y no un sistema de diagnóstico:
sus resultados no reemplazan la evaluación de un especialista.
\pendiente{evidencia local, si el equipo la incorpora: acceso a TC o a
radiólogos en el Perú}

\subsection{Usuarios y procesos beneficiados}\label{sec:users}
Thoraxis está pensado para usuarios del ámbito de la investigación, en cuatro
procesos:

\begin{itemize}
    \item \textbf{Equipos de investigación en imagen médica} que analizan
          cohortes retrospectivas de TC y necesitan medir el volumen de cada
          pulmón en muchos estudios, sin delimitarlos uno por uno.
    \item \textbf{Radiólogos y anotadores} que preparan conjuntos de datos: en
          lugar de trazar el pulmón desde cero, una tarea que puede tomar horas
          por estudio \parencite{ma2021}, revisan y corrigen la máscara que
          propone el modelo.
    \item \textbf{Desarrolladores de otros modelos}, como detectores de nódulos
          o de lesiones, que necesitan el pulmón delimitado como paso previo
          de su análisis \parencite{hofmanninger2020}.
    \item \textbf{Desarrolladores de sistemas} que quieren integrar la
          segmentación en sus aplicaciones: el modelo se distribuye en ONNX y
          se expone mediante una API, de modo que puede ejecutarse desde otros
          lenguajes y marcos.
\end{itemize}

\noindent En todos los casos, el resultado lo revisa una persona. Thoraxis no
está dirigido a la atención de pacientes: usarlo en la práctica clínica
exigiría una validación clínica y las autorizaciones sanitarias
correspondientes, que quedan fuera del alcance del proyecto.

\subsection{Objetivos}\label{sec:objectives}

\paragraph{Objetivo general.}
Desarrollar y evaluar de forma reproducible un modelo de Deep Learning que
segmente automáticamente el pulmón izquierdo y el derecho en tomografías
computarizadas de tórax, incluidos pulmones con patología, e integrarlo en un
prototipo que muestre la segmentación y calcule el volumen de cada pulmón.

\paragraph{Objetivos específicos.}
\begin{enumerate}
    \item Preparar un conjunto de TC de tórax de fuentes públicas, con máscaras
          de referencia de ambos pulmones, un preprocesamiento homogéneo y una
          partición por paciente sin fugas entre entrenamiento, validación y
          prueba.
    \item Implementar una solución base, una U-Net 2D entrenada corte por
          corte, y una propuesta que la mejore con contexto entre cortes
          contiguos, aumentos de datos y datos de entrenamiento más diversos.
    \item Evaluar ambas soluciones en las mismas condiciones, con el
          coeficiente de Dice y la distancia de Hausdorff al percentil 95
          (HD95) e intervalos de confianza por \textit{bootstrap}, en una
          prueba interna y en una validación externa con datos de otras
          instituciones, usando TotalSegmentator como referencia publicada.
    \item Analizar los errores de la propuesta, en especial en pulmones con
          patología, identificar sus limitaciones y estimar la incertidumbre de
          cada segmentación para explicar dónde y cuándo el modelo es menos
          confiable.
    \item Exportar el modelo a ONNX y desplegarlo en un servicio de inferencia
          y un visor web que muestren la máscara sobre los cortes y el volumen
          de cada pulmón.
    \item Registrar los experimentos, configuraciones y versiones con MLflow,
          de modo que la evaluación pueda reproducirse desde los
          repositorios.
\end{enumerate}

\subsection{Alcance}\label{sec:scope}
El proyecto incluye:

\begin{itemize}
    \item Tomografías computarizadas de tórax de conjuntos de datos públicos y
          anonimizados, descritos en la Sección~\ref{sec:data}.
    \item La segmentación del pulmón izquierdo y del derecho, con modelos que
          procesan cortes individuales (2D) o pilas de cortes contiguos
          (2.5D).
    \item El cálculo del volumen de cada pulmón a partir de la máscara y del
          tamaño del vóxel.
    \item Un mapa de incertidumbre por vóxel y una puntuación de confianza por
          estudio, como forma de explicabilidad.
    \item Un prototipo con un servicio de inferencia y un visor web para
          cargar un estudio de prueba y revisar el resultado.
\end{itemize}

\noindent El proyecto no incluye:

\begin{itemize}
    \item La segmentación de lesiones, como la infección por COVID-19 o los
          nódulos, ni de los lóbulos u otras estructuras del tórax, que quedan
          como extensiones futuras.
    \item El diagnóstico, la clasificación de enfermedades o cualquier
          recomendación clínica.
    \item El uso con datos de pacientes fuera de los conjuntos públicos, la
          validación clínica prospectiva y la certificación como dispositivo
          médico.
    \item Otras modalidades de imagen, como la radiografía de tórax o la
          resonancia magnética.
\end{itemize}

\subsection{Resultado esperado y criterio de éxito}\label{sec:expected-outcome}
El resultado esperado es un modelo de segmentación de ambos pulmones,
exportado a ONNX y versionado, un prototipo que lo usa, y el código, las
configuraciones y los registros necesarios para reproducir su evaluación. Se
considerará alcanzado si se cumplen los siguientes criterios:

\begin{enumerate}
    \item \textbf{Exactitud externa:} el Dice medio por paciente es de al
          menos 0,95 para cada pulmón en la validación externa, reportado con
          su intervalo de confianza del 95~\%.
    \item \textbf{Mejora frente a la base:} la propuesta supera a la solución
          base en Dice en la prueba interna y en la externa, y el intervalo de
          confianza de la diferencia por pares no incluye el cero. La HD95 se
          reporta y se compara de la misma forma.
    \item \textbf{Prototipo funcional:} el prototipo procesa un estudio de
          prueba de extremo a extremo y muestra la máscara y el volumen de cada
          pulmón.
    \item \textbf{Reproducibilidad:} otro integrante del equipo reproduce la
          evaluación desde los repositorios y obtiene las mismas métricas.
\end{enumerate}
